% Fragment partage : inclus par abonnement.tex (dossier autonome)
% et par dossier-conception-alanya.tex (dossier client assemble).
% Ne pas y mettre de preambule ni de \part : c'est l'appelant qui les pose.

\chapter{L'offre}

\section{Une offre, deux durées}

Il n'y a qu'une offre, \textbf{Alanya Plus}, vendue sous
deux durées~:

\begin{center}
\begin{tabularx}{\linewidth}{@{}p{3.2cm}p{3.4cm}X@{}}
\toprule
\thd{Durée} & \thd{Prix} & \thd{Relance avant l'échéance} \\
\midrule
Mensuelle & 250~FCFA & 7~jours avant \\
Annuelle  & 2\,000~FCFA & 30~jours avant --- soit 167~F par mois, quatre mois offerts \\
\bottomrule
\end{tabularx}
\end{center}

Les deux durées donnent exactement la même chose. L'annuelle est mise en avant dans l'écran
d'offre~: elle est moins chère au mois et elle réduit le nombre de renouvellements --- donc
d'occasions de perdre un abonné pour un simple oubli.

\section{Ce qui se paie}

\begin{center}
\begin{tabularx}{\linewidth}{@{}p{4.4cm}X@{}}
\toprule
\thd{Fonctionnalité} & \thd{Ce qu'elle apporte} \\
\midrule
Traduction des messages & Traduire un message, ou toute une conversation automatiquement,
sur le téléphone. \\
Sauvegarde des données & Sauvegardes chiffrées régulières, sur le téléphone ou sur Google
Drive. \\
Trajets de confiance & Partager sa route à sa liste \emph{Confiance} et être attendu à
l'arrivée. \\
Sonneries par liste & Une sonnerie propre à la famille, une autre au bureau. \\
Personnalisation du style & \emph{Plus tard.} Couleurs et fonds au-delà du clair et du
sombre. \\
La coche \badgePastilleCoche[1] & Le \emph{droit} à la coche --- qui demande en plus une
vérification d'identité (chapitre~\ref{chap:abo-coche}). \\
\bottomrule
\end{tabularx}
\end{center}

\section{Ce qui reste gratuit, toujours}

Écrire, appeler, publier un statut, créer un groupe ou une réunion, envoyer des médias~: tout
ce qui fait d'Alanya une messagerie reste gratuit. On ne fait pas payer la conversation.

Trois cas méritent d'être dits explicitement, parce qu'on pourrait les croire payants~:

\begin{center}
\begin{tabularx}{\linewidth}{@{}p{4.8cm}X@{}}
\toprule
\thd{Reste gratuit} & \thd{Pourquoi} \\
\midrule
\emph{Suivre} le trajet d'un proche & Seul celui qui \emph{partage} sa route est abonné. Ses
cinq proches ne doivent rien payer pour veiller sur lui. \\
Restaurer une sauvegarde & Faire une \emph{nouvelle} sauvegarde est payant~; récupérer ses
propres données ne l'est jamais. \\
Le thème clair ou sombre & Il existe déjà, et il relève de l'accessibilité autant que du
goût. \\
\bottomrule
\end{tabularx}
\end{center}

\begin{note}[Le SOS fait partie de l'offre]
Le bouton SOS est réservé aux abonnés, comme les trajets de confiance dont il fait partie.
Une exception~: un trajet lancé pendant l'abonnement va à son terme avec son SOS, même si
l'échéance tombe en route. On ne coupe pas une protection en cours de chemin.
\end{note}

\section{Des plans qui se règlent sans développeur}

Prix, durée, délai de relance, fonctionnalités incluses, mise en avant, disponibilité~: tout
se modifie depuis l'administration. Deux règles encadrent ces modifications~:

\begin{itemize}[leftmargin=1.4em]
  \item \textbf{Un changement de prix ne touche jamais une période déjà payée.} Il s'applique
        au prochain renouvellement, et l'abonné en est informé dans sa relance.
  \item \textbf{Un plan ne se supprime pas, il se retire.} Il disparaît de l'offre, mais
        l'historique de ceux qui l'ont payé reste lisible.
\end{itemize}

\begin{note}[Ce que l'administration ne peut pas faire seule]
Elle peut rendre une fonctionnalité existante payante ou gratuite, et choisir ce que chaque
plan inclut. Elle ne peut pas \emph{inventer} une fonctionnalité~: le verrou doit exister
dans le code. C'est pourquoi la personnalisation du style figure déjà au catalogue, marquée
\guillemotleft~à venir~\guillemotright{} --- elle s'activera le jour où l'application saura
l'appliquer.
\end{note}

% =====================================================================
\chapter{La coche indigo}
\label{chap:abo-coche}

\section{Deux conditions, pas une}

La coche \badgePastilleCoche[1] s'affiche quand \textbf{les deux} conditions sont réunies~:

\begin{center}
\begin{tabularx}{\linewidth}{@{}p{3.6cm}p{4.6cm}X@{}}
\toprule
\thd{Identité} & \thd{Abonnement} & \thd{Ce qui s'affiche} \\
\midrule
non vérifiée    & actif ou non   & aucune coche --- l'abonné a ses fonctionnalités \\
en cours d'examen & actif        & aucune coche --- \guillemotleft~en cours~\guillemotright{}
dans son profil \\
\textbf{vérifiée} & \textbf{actif} & \badgePastilleCoche[1]~\textbf{la coche} \\
vérifiée        & échu           & aucune coche --- elle revient dès le réabonnement \\
révoquée        & actif ou non   & aucune coche, jusqu'à nouvel examen \\
\bottomrule
\end{tabularx}
\end{center}

L'abonnement seul ne suffit pas~: sans contrôle, n'importe qui pourrait acheter pour 250~F une
coche au nom d'un autre. L'identité seule ne suffit plus~: hors phase gratuite, la coche fait
partie de l'offre.

\section{Ce qu'elle atteste}

\guillemotleft~Alanya a vérifié que la personne derrière ce compte est bien celle dont il porte
le nom.~\guillemotright{} Rien de plus. Ce n'est ni une recommandation, ni une marque de
notoriété.

\begin{alerte}[Le sens de la coche change par rapport au volet 5]
Le volet~5 réservait la coche personnelle aux \emph{personnalités publiques}, sur preuve de
notoriété. Elle devient accessible à \emph{tout abonné} dont l'identité est vérifiée. La pièce
de notoriété disparaît donc des pièces demandées~; restent la pièce d'identité et un selfie.
\end{alerte}

\section{Quand elle tombe}

\begin{itemize}[leftmargin=1.4em]
  \item \textbf{À l'échéance}, si l'abonnement n'est pas renouvelé. L'identité, elle, reste
        vérifiée~: se réabonner fait revenir la coche immédiatement, sans nouveau dossier.
  \item \textbf{Sur révocation} par un administrateur, avec un motif --- usurpation, fraude.
        C'est la seule chute qui n'a rien à voir avec l'argent.
  \item \textbf{Quand le compte change de nom.} La vérification portait sur un nom~; le
        changer retire la coche jusqu'à ce qu'un administrateur ait revu le dossier. Sans
        cette règle, on achète une coche sous son vrai nom, puis on se renomme
        \guillemotleft~Orange Cameroun~\guillemotright.
\end{itemize}

% =====================================================================
\chapter{Le parcours de l'utilisateur}

\section{Le profil met l'offre en avant}

Une carte se place \textbf{juste sous l'en-tête du profil}, avant les contacts préférés. C'est
l'emplacement le plus vu de l'écran~; son contenu change avec la situation de l'utilisateur.

\begin{center}
  \imageOuCadre{maquettes/abonnement-profil.png}{\linewidth}%
    {La carte Plus du profil dans ses cinq états\\
     \texttt{maquettes/abonnement-profil.png}}
\end{center}

\begin{center}
\begin{tabularx}{\linewidth}{@{}p{4.4cm}X@{}}
\toprule
\thd{Situation} & \thd{Ce que dit la carte} \\
\midrule
Phase gratuite & \guillemotleft~Obtenez la coche --- gratuite pendant le lancement.~\guillemotright
Il n'y a rien à vendre~; on invite à la vérification. \\
Période de grâce & \guillemotleft~Les fonctionnalités Plus deviennent payantes le 12
octobre.~\guillemotright{} Avec le prix, et un bouton pour s'abonner dès maintenant sans
perdre les jours gratuits restants. \\
Non abonné & \guillemotleft~Passer à Alanya Plus~\guillemotright{}~: les cinq icônes des
fonctionnalités, \guillemotleft~dès 167~F par mois~\guillemotright, un bouton. \\
Abonné & L'échéance, le renouvellement automatique, et --- si l'identité n'est pas encore
vérifiée --- l'invitation à obtenir la coche. \\
Échéance proche ou dépassée & Le compte à rebours, et ce qui sera supprimé si rien n'est
fait. \\
\bottomrule
\end{tabularx}
\end{center}

Une entrée \guillemotleft~Mon abonnement~\guillemotright{} rejoint aussi le menu du profil~: la
carte attire, l'entrée permet de retrouver ses paiements et de régler le renouvellement.

\begin{note}[Maquette interactive, mobile]
\href{https://claude.ai/code/artifact/6ac3fb3e-7f18-4e2d-b07c-c08127ef61e4}%
{\texttt{claude.ai/code/artifact/6ac3fb3e}}\\
Source versionnée~: \code{docs/architecture/maquettes/abonnement.html}~; captures produites
par \code{maquettes/\_capture\_abonnement.py}.
\end{note}

\section{S'abonner}

\begin{center}
  \imageOuCadre{maquettes/abonnement-parcours.png}{\linewidth}%
    {Offre, paiement mobile money, confirmation\\
     \texttt{maquettes/abonnement-parcours.png}}
\end{center}

Trois écrans~: l'\textbf{offre} (les deux durées, l'annuelle présélectionnée, ce qui est
inclus), le \textbf{paiement} (Orange Money ou MTN MoMo, le numéro à débiter), puis
l'\textbf{attente de confirmation} --- l'utilisateur valide sur son téléphone avec son code
secret, l'écran se met à jour de lui-même.

\begin{note}[Pourquoi l'écran attend au lieu de conclure]
Un paiement mobile money n'est jamais confirmé au moment où on le demande~: il l'est quand
l'opérateur prévient Alanya, quelques secondes ou minutes plus tard. L'écran le montre tel
quel plutôt que d'afficher un succès qui pourrait se révéler faux. Le paiement simulé imite
ce délai, pour que l'écran soit éprouvé avant l'arrivée du vrai agrégateur.
\end{note}

\section{Rencontrer une fonctionnalité verrouillée}

Toucher une fonctionnalité payante sans abonnement n'affiche pas d'erreur~: un panneau monte
du bas de l'écran, dit ce que la fonctionnalité apporte, et propose l'offre. Refermer le
panneau ramène exactement où l'on était.

\section{L'échéance}

\begin{center}
  \imageOuCadre{maquettes/abonnement-echeance.png}{\linewidth}%
    {Relance, panneau de fonctionnalité verrouillée, abonnement expiré\\
     \texttt{maquettes/abonnement-echeance.png}}
\end{center}

\begin{center}
\begin{tikzpicture}[
  x=1cm, y=1cm,
  j/.style={font=\sffamily\scriptsize, text=AlanyaGray, align=center},
  e/.style={font=\sffamily\small\bfseries, text=AlanyaIndigoDark, align=center}
]
  \fill[AlanyaSuccess!14] (0,0) rectangle (6.2,0.5);
  \fill[warning!18]       (6.2,0) rectangle (9.4,0.5);
  \fill[danger!10]        (9.4,0) rectangle (15.4,0.5);
  \draw[AlanyaLine] (0,0) rectangle (15.4,0.5);
  \foreach \x/\t in {6.2/Relance, 9.4/Échéance, 15.4/Suppression}{
    \draw[AlanyaGray, semithick] (\x,-0.12) -- (\x,0.62);
    \node[e] at (\x,1.05) {\t};
  }
  \draw[AlanyaGray, semithick, dashed] (13.0,-0.12) -- (13.0,0.62);
  \node[j] at (3.1,0.25)  {abonné};
  \node[j] at (7.8,0.25)  {J-7 ou J-30};
  \node[j] at (11.2,0.25) {données conservées};
  \node[j] at (6.2,-0.45) {\code{reminder\_days}};
  \node[j] at (9.4,-0.45) {coche et fonctions coupées};
  \node[j] at (13.0,-0.45) {dernier avis, J-7};
  \node[j, anchor=east] at (15.4,-0.45) {N jours après};
\end{tikzpicture}
\end{center}

\begin{itemize}[leftmargin=1.4em]
  \item \textbf{La relance} arrive 7~jours avant la fin d'un abonnement mensuel, 30~jours
        avant celle d'un annuel --- par notification et par la carte du profil.
  \item \textbf{À l'échéance}, les fonctionnalités payantes s'arrêtent et la coche tombe.
        Rien n'est encore effacé.
  \item \textbf{Pendant N~jours} (réglables, 30 par défaut), les données restent en place~:
        se réabonner rend tout, exactement comme avant.
  \item \textbf{Une semaine avant la suppression}, un dernier avertissement.
  \item \textbf{Au bout de N~jours}, les données propres aux fonctionnalités payantes sont
        supprimées --- le détail, fonctionnalité par fonctionnalité, est en partie~II.
\end{itemize}

\section{Le renouvellement automatique}

L'utilisateur peut l'activer à tout moment dans \guillemotleft~Mon abonnement~\guillemotright.
Ce qu'il fait réellement dépend du moyen de paiement~:

\begin{center}
\begin{tabularx}{\linewidth}{@{}p{3.8cm}X@{}}
\toprule
\thd{Moyen} & \thd{Ce qui se passe à l'échéance} \\
\midrule
Paiement simulé & Le renouvellement réussit tout seul --- c'est ce qui permet de tester le
parcours complet dès maintenant. \\
Mobile money & L'opérateur envoie une demande de paiement la veille~; l'utilisateur la valide
avec son code. Un vrai prélèvement sans action n'existe pas sur mobile money. \\
Porte-monnaie Alanya & \emph{Plus tard.} Débit direct du solde, sans action --- le seul vrai
renouvellement automatique hors stores. \\
App Store, Play Store & Le store renouvelle lui-même, selon ses règles. \\
\bottomrule
\end{tabularx}
\end{center}

\section{Obtenir la coche}

\begin{center}
  \imageOuCadre{maquettes/abonnement-coche.png}{\linewidth}%
    {Pièce d'identité, selfie, suivi du dossier\\
     \texttt{maquettes/abonnement-coche.png}}
\end{center}

Le parcours est celui du volet~5~: photographier sa pièce d'identité, prendre un selfie en la
tenant, envoyer. L'écran annonce \emph{avant} l'envoi que les pièces sont chiffrées,
consultées seulement par l'administration, et supprimées 90~jours après la décision. Réponse
annoncée \guillemotleft~sous 5~jours ouvrés~\guillemotright.

% =====================================================================
\chapter{Du gratuit au payant}

\section{Trois phases}

L'administration décide quand l'offre devient payante. Jusque-là, tout est gratuit~; ensuite,
chacun dispose d'une période de grâce avant que les fonctionnalités ne soient réservées aux
abonnés.

\begin{center}
\begin{tikzpicture}[
  x=1cm, y=1cm,
  p/.style={minimum height=0.9cm, font=\sffamily\small, text=AlanyaInk, align=center,
            inner sep=0pt},
  m/.style={font=\sffamily\scriptsize, text=AlanyaGray, align=center}
]
  \node[p, fill=AlanyaSuccess!14, minimum width=5.6cm] at (2.8,0)
    {\textbf{Gratuit}\\[-1pt]{\scriptsize tout pour tous}};
  \node[p, fill=warning!18, minimum width=4.4cm] at (7.8,0)
    {\textbf{Grâce}\\[-1pt]{\scriptsize 30 jours par défaut}};
  \node[p, fill=AlanyaContainer, minimum width=5.4cm] at (12.7,0)
    {\textbf{Payant}\\[-1pt]{\scriptsize réservé aux abonnés}};
  \draw[AlanyaIndigo, very thick] (5.6,-0.62) -- (5.6,0.62);
  \draw[AlanyaIndigo, very thick] (10.0,-0.62) -- (10.0,0.62);
  \node[m] at (5.6,-1.0)  {l'admin active};
  \node[m] at (10.0,-1.0) {fin de la grâce};
\end{tikzpicture}
\end{center}

\begin{center}
\begin{tabularx}{\linewidth}{@{}p{2.2cm}XX@{}}
\toprule
 & \thd{Fonctionnalités} & \thd{Coche} \\
\midrule
Gratuit & Pour tous. L'offre n'est pas proposée à l'achat. & Sur vérification d'identité
seule. \\
Grâce & Pour tous. On peut s'abonner~; la première période commence \emph{à la fin} de la
grâce, les jours gratuits ne sont pas perdus. & Conservée par ceux qui l'ont. Elle tombe à la
fin de la grâce s'ils ne s'abonnent pas. \\
Payant & Pour les abonnés, et pendant l'éventuel essai gratuit des nouveaux inscrits. &
Identité vérifiée \emph{et} abonnement actif. \\
\bottomrule
\end{tabularx}
\end{center}

\begin{note}[Un seul bouton, pas trois]
L'administrateur ne choisit pas une phase~: il n'a qu'une action, \guillemotleft~Activer le
payant~\guillemotright. Elle déclenche la grâce, qui débouche d'elle-même sur la phase payante
à sa date de fin. Il n'existe pas de bouton qui rende l'offre payante sur-le-champ~: la grâce
dure au moins 7~jours, pour que personne ne perde un accès sans avoir été prévenu. Le bandeau
des trois phases, en tête de la page d'administration, est un repère, pas un sélecteur.
\end{note}

\section{Ce que voit l'utilisateur au passage}

\begin{itemize}[leftmargin=1.4em]
  \item \textbf{À l'activation}~: un message du compte officiel à tous les comptes, qui
        annonce la date et le prix. La carte du profil passe en mode
        \guillemotleft~grâce~\guillemotright.
  \item \textbf{Sept jours avant la fin de la grâce}~: un second message, aux seuls non-abonnés.
  \item \textbf{À la fin de la grâce}~: les fonctionnalités se ferment pour les non-abonnés,
        avec le même délai de N~jours avant suppression qu'une échéance ordinaire.
\end{itemize}

\section{Les réglages de l'interrupteur}

\begin{center}
\begin{tabularx}{\linewidth}{@{}p{4.6cm}p{2.2cm}X@{}}
\toprule
\thd{Réglage} & \thd{Défaut} & \thd{Effet} \\
\midrule
Payant activé & non & L'interrupteur lui-même. \\
Durée de la grâce & 30~jours & 7~jours au minimum. Posée à l'activation~; la date de fin
reste modifiable ensuite, pour prolonger. \\
Essai gratuit des nouveaux inscrits & 0~jour & Après la grâce, un inscrit reçoit cette durée
d'accès complet. \\
Conservation après échéance & 30~jours & Le \guillemotleft~N~\guillemotright{} avant
suppression des données payantes. \\
\bottomrule
\end{tabularx}
\end{center}

\begin{alerte}[Pas d'activation tant que le paiement est simulé]
En production, l'interrupteur refuse de s'activer tant que le fournisseur de paiement est le
simulateur. Sinon, n'importe qui obtiendrait un abonnement réel avec un paiement fictif.
\end{alerte}

% =====================================================================
\chapter{Le parcours de l'administrateur}

\section{L'interrupteur et ses réglages}

\begin{center}
  \imageOuCadre{maquettes/abonnement-admin-reglages.png}{\linewidth}%
    {Phase courante, interrupteur, grâce, essai, conservation\\
     \texttt{maquettes/abonnement-admin-reglages.png}}
\end{center}

La page dit d'abord dans quelle phase on est, et ce que l'activation déclencherait~: combien
de comptes recevront l'annonce, combien de coches entreront en grâce. Activer demande une
confirmation qui rappelle la date de fin de grâce.

\begin{note}[Maquette interactive, administration]
\href{https://claude.ai/code/artifact/ffa8c829-4047-4f86-b089-3330fe810687}%
{\texttt{claude.ai/code/artifact/ffa8c829}}\\
Source versionnée~: \code{docs/architecture/maquettes/abonnement-admin.html}.
\end{note}

\section{Les plans}

\begin{center}
  \imageOuCadre{maquettes/abonnement-admin-plans.png}{\linewidth}%
    {Liste des plans et éditeur\\
     \texttt{maquettes/abonnement-admin-plans.png}}
\end{center}

Pour chaque plan~: son nom dans les trois langues, sa durée, son prix, son délai de relance,
les fonctionnalités incluses (cases à cocher), s'il est mis en avant, s'il est proposé. Un
aperçu montre la carte telle qu'elle apparaîtra dans l'application.

\section{Un utilisateur}

\begin{center}
  \imageOuCadre{maquettes/abonnement-admin-fiche.png}{\linewidth}%
    {Abonnement, coche et paiements sur la fiche utilisateur\\
     \texttt{maquettes/abonnement-admin-fiche.png}}
\end{center}

La fiche utilisateur gagne une carte \guillemotleft~Abonnement et coche~\guillemotright~:
la période en cours, l'historique des périodes et des paiements, l'état du dossier
d'identité. Deux actions~: \textbf{offrir} un abonnement (durée et motif obligatoires --- un
geste commercial, un partenaire, un testeur) et \textbf{révoquer} la coche (motif
obligatoire). Le sélecteur \guillemotleft~État de vérification~\guillemotright{} de la carte
\guillemotleft~Socle~\guillemotright{} disparaît~: l'état se déduit des décisions, il ne se
saisit plus.

\section{Les dossiers d'identité}

L'écran d'instruction est celui du volet~5~: ce que le compte déclare à gauche, les pièces à
droite, trois issues --- approuver, refuser avec motif, demander une pièce. Une file dédiée
les rassemble, avec un compteur dans le menu.

% =====================================================================
\chapter{Le paiement~: simulé aujourd'hui, réel demain}

\section{Un faux agrégateur, un vrai parcours}

Tant qu'aucun agrégateur n'est choisi, le serveur en \emph{imite} un~: il reçoit la demande,
fait patienter quelques secondes, puis répond --- le plus souvent par un succès. Certains
numéros de test provoquent volontairement un échec, un refus ou une absence de réponse, pour
que chaque cas soit vu à l'écran avant d'arriver en vrai.

Le jour où l'agrégateur est choisi, on écrit le module qui lui parle et on le met à la place
du simulateur. L'application, l'administration, les abonnements et la coche ne changent pas.

\section{Les stores}

\begin{alerte}[Apple et Google imposent leurs règles]
L'application sera publiée sur l'App Store et le Play Store. Pour débloquer des
fonctionnalités numériques, Apple impose son propre système de paiement, et Google le sien
dans la plupart des pays --- avec une commission de 15 à 30\,\%. Mobile money et, plus tard,
porte-monnaie ne seront donc pas forcément proposables partout dans l'application.

La conception l'absorbe~: un store est un fournisseur de paiement de plus, au même endroit que
le simulateur. Mais la règle exacte pour le Cameroun, et le prix affiché sur chaque store,
restent \textbf{à décider avant de brancher un paiement réel}.
\end{alerte}
